\documentclass[sigconf,review,anonymous]{acmart}

\AtBeginDocument{%
  \providecommand\BibTeX{{%
    Bib\TeX}}}

% RecSys 2026 submission mode (anonymous review).
% For camera-ready, switch documentclass to [sigconf] and replace rights metadata
% below with the exact values from ACM rights management.
\setcopyright{none}
\acmConference[RecSys '26]{Proceedings of the 20th ACM Conference on Recommender Systems}{September 28--October 2, 2026}{Minneapolis, MN, USA}
\acmBooktitle{Proceedings of the 20th ACM Conference on Recommender Systems (RecSys '26), September 28--October 2, 2026, Minneapolis, MN, USA}
\acmDOI{}
\acmISBN{}

\usepackage{array}
\usepackage{multirow}
\usepackage{tabularx}

\newcolumntype{Y}{>{\centering\arraybackslash}X}
\newcommand{\tblbest}[1]{\textbf{#1}}
\newcommand{\tblsecond}[1]{\underline{#1}}
\newcommand{\placeholderpanel}[1]{\fbox{\rule{0pt}{#1}\rule{0.92\linewidth}{0pt}}}

\begin{document}

\title[RouteRec for Sequential Recommendation]{RouteRec: Behavior-Guided Expert Routing for Sequential Recommendation}

\author{First Author}
\affiliation{%
  \institution{KAIST}
  \department{Kim Jaechul Graduate School of AI}
  \city{Daejeon}
  \country{Republic of Korea}}
\email{first.author@kaist.ac.kr}

\author{Second Author}
\affiliation{%
  \institution{KAIST}
  \department{Kim Jaechul Graduate School of AI}
  \city{Daejeon}
  \country{Republic of Korea}}
\email{second.author@kaist.ac.kr}

\renewcommand{\shortauthors}{Author et al.}

\begin{abstract}
Sequential recommendation exhibits heterogeneous behavioral regimes, yet most models still process all prefixes through largely shared computation. In practice, some prefixes are governed by stable preference, others by evolving session intent, and others by immediate local transitions. This creates a mismatch between behavioral diversity and shared computation. Mixture-of-experts (MoE) offers a natural mechanism for conditional computation, but hidden-state-driven routing remains limited for this setting: it is often semantically opaque and does not explicitly align expert allocation with different temporal scopes.

We propose \textbf{RouteRec}, a behavior-guided MoE framework for sequential recommendation built around three principles. First, RouteRec derives \emph{portable behavioral cues} from ordinary interaction logs, avoiding dependence on heavy side-information encoders. Second, it applies \emph{coarse-to-fine multi-stage routing} across macro, mid, and micro scopes so that long-range tendency, within-session evolution, and immediate transition cues can influence routing at different granularities. Third, it uses \emph{behavior-guided expert allocation} that decouples routing signals from hidden representations, making expert selection more inspectable and behavior-grounded.

This design treats behavior features as routing controls rather than as yet another representation-fusion channel. As a result, RouteRec preserves strong sequence modeling in the representation path while introducing a lightweight and transferable control interface for conditional computation. Experiments show that RouteRec improves ranking quality on several datasets while remaining competitive on others, and yields more behavior-consistent routing patterns in diagnostic analyses.
\end{abstract}

\begin{CCSXML}
<ccs2012>
 <concept>
  <concept_id>10002951.10003317.10003347.10003350</concept_id>
  <concept_desc>Information systems~Recommender systems</concept_desc>
  <concept_significance>500</concept_significance>
 </concept>
</ccs2012>
\end{CCSXML}

\ccsdesc[500]{Information systems~Recommender systems}

\keywords{sequential recommendation, mixture of experts, conditional computation, routing, behavioral features}

\maketitle

\section{Introduction}
Sequential recommendation is usually formulated as next-item prediction over an ordered interaction history, but the underlying behavior is rarely homogeneous. Even for the same user, one prefix may reflect stable preference, another may reflect an evolving within-session goal, and another may be dominated by immediate local exploration. Prior work on sequential and session-based recommendation has repeatedly shown that user behavior mixes multiple regimes rather than following a single stationary pattern \cite{wang2019seqsurvey,wang2021sessionsurvey}.

Strong sequential recommenders model these histories increasingly well \cite{hidasi2016gru4rec,kang2018sasrec,li2020tisasrec,qiu2022duorec,shin2024bsarec,liu2025sigma}. However, most of them still rely on largely shared computation across prefixes that may arise from very different behavioral contexts. This shared-path assumption is practical and often effective on average, but it also creates a modeling mismatch: heterogeneous behavior is acknowledged in the data, while computation remains mostly uniform.

Figure~\ref{fig:intro-behavioral-features} illustrates this gap. Distinct users and sessions can differ in pace, repetition, focus, and local transition patterns, yet standard sequential models expose limited control over how such differences should change the computation path. Existing advances therefore improve representation quality, but they leave the allocation of computation mostly implicit.

\begin{figure}[!t]
  \centering
  \includegraphics[width=\columnwidth]{figures/figure1_different_user.pdf}
  \caption{Different users and sessions can require different computational paths.}
  \Description{A problem-motivation figure contrasting different user and session behavior patterns, highlighting that behaviorally different cases are often processed with the same computation path in standard sequential recommendation.}
  \label{fig:intro-behavioral-features}
\end{figure}

Feature-aware and heterogeneity-aware sequential models partially address this issue, but mostly by enriching representations, reweighting attention, or decomposing interests inside a shared computation pipeline. In contrast, the question we ask here is more direct: can observable behavior cues control \emph{which computation path} should be used? Mixture-of-experts (MoE) is a natural candidate for this purpose \cite{jacobs1991adaptive,jordan1994hme,shazeer2017moe,fedus2022switch}, because it can replace one shared transformation with condition-dependent expert allocation. Yet simply inserting experts is not enough. For sequential recommendation, routing must still answer why specialization is needed in the first place, why routing should depend on temporal scope rather than a single view, and why the routing signal should remain interpretable instead of being absorbed entirely into hidden states.

\textbf{Challenge 1: behavioral diversity versus shared computation.} Real prefixes arise from different behavioral regimes, but most sequential recommenders still process them with largely shared computation. This leaves the model with limited ability to specialize computation for persistent preference, evolving session intent, or immediate local transitions. To address this, we need routing cues that are portable and derivable from ordinary logs rather than from heavy dataset-specific side branches.

\textbf{Challenge 2: temporal-scope mismatch.} Once routing is introduced, another issue appears immediately: behavioral evidence arrives at different horizons. Long-range history reflects relatively stable tendency, within-session context reflects evolving intent, and short-range context reflects immediate transition pressure. A single routing view is therefore too coarse: it can blur signals that should influence expert allocation at different granularities. To address this, routing should be aligned explicitly with macro, mid, and micro temporal scopes.

\textbf{Challenge 3: opaque hidden-state routing.} Even if multi-stage routing is available, routing is often driven primarily by hidden representations. This is effective, but it makes expert allocation difficult to inspect because the same latent space must both encode sequence content and decide computation paths. If routing is meant to reflect behavior, it should be grounded in signals that can be traced back to observable interaction patterns rather than hidden-state geometry alone.

We address these challenges with \textbf{RouteRec}, a behavior-guided MoE framework for sequential recommendation. RouteRec is organized around three design axes that map directly to the challenges above. \emph{Portable Behavioral Cues} summarize ordinary logs into compact routing signals, providing a lightweight interface for behavior-aware control. Such cues include, for example, inter-event tempo, category concentration and switching, recurrence patterns, and popularity exposure. \emph{Coarse-to-Fine Multi-Stage Routing} applies those cues across macro, mid, and micro scopes so that expert allocation follows temporal evidence granularity rather than a single global bias. \emph{Behavior-Guided Expert Allocation} decouples routing signals from hidden representations, so experts transform sequence states while routing remains tied to inspectable behavior cues.

This framing deliberately treats behavioral information as a control interface for conditional computation rather than as another feature-fusion pathway. As a result, RouteRec preserves a strong sequence backbone while exposing a lightweight mechanism for behavior-grounded specialization. The rest of the paper follows this structure: Related Work clarifies where prior studies stop short of these goals, Method formalizes the three design axes, and Experiments test whether compact cues, multi-stage routing, and decoupled expert allocation each matter in practice.

The main contributions of this work are as follows:
\begin{itemize}
  \item We identify three front-end challenges for conditional computation in sequential recommendation: the mismatch between behavioral diversity and shared computation, temporal-scope mismatch in routing, and the opacity of hidden-state-only expert allocation.
  \item We propose RouteRec, which combines portable behavioral cues, coarse-to-fine multi-stage routing, and behavior-guided expert allocation to align expert usage with observable behavior.
  \item We design RouteRec to operate from ordinary interaction logs, making the routing interface lightweight, transferable, and suitable for diagnostic analysis across datasets.
  \item We show that RouteRec improves recommendation quality over strong baselines on several datasets and yields more behavior-consistent routing patterns in ablation and diagnostic studies.
\end{itemize}

\section{Related Work}
\label{sec:related}

\subsection{Shared-Path Sequential Modeling and Conditional Computation}
Sequential recommendation has progressed from recurrent and convolutional models to graph and Transformer architectures \cite{hidasi2016gru4rec,tang2018caser,wu2019srgnn,kang2018sasrec,sun2019bert4rec}. Recent methods further improve performance with time-aware attention, contrastive learning, and structural inductive bias \cite{li2020tisasrec,xie2020cl4srec,wang2023multiplecontrast,qiu2022duorec,du2023fearec,shin2024bsarec,liu2025sigma}. These methods substantially strengthen sequence representation quality, but they generally preserve one shared computation pathway for behaviorally different prefixes.

Conditional computation and MoE provide a natural response to this limitation. Classical MoE and later sparse routing work show that different inputs need not share the same transformation path \cite{jacobs1991adaptive,jordan1994hme,shazeer2017moe,fedus2022switch}. Recommendation systems have also adopted expert decomposition to improve specialization and capacity \cite{ma2018mmoe,tang2020ple,liu2025fame,bian2023m3srec,zhang2025hm4sr}. However, existing sequential recommendation literature still gives limited attention to the specific question we focus on here: how to let observable behavioral diversity determine computation allocation. RouteRec addresses this gap by building a portable behavior cue interface for routing rather than keeping all heterogeneity inside a shared-path encoder.

\subsection{Multi-Scope Behavioral Signals in Sequential Recommendation}
Many studies explicitly acknowledge that user state is heterogeneous or non-stationary, including short-term intent modeling, multi-interest extraction, and long-short preference decomposition \cite{li2017narm,liu2018stamp,ying2018shan,ma2019hgn,li2019mind,cen2020comirec,quadrana2017hrnn}. Feature-aware sequential recommenders also incorporate timestamps, side information, and auxiliary feature channels to enrich sequence encoding \cite{li2020tisasrec,zhang2019fdsa,xie2022difsr,du2023fearec,bauer2024hybridsession}. Together, these lines of work show that sequential behavior is expressed across multiple temporal scopes and can benefit from richer behavioral evidence.

At the same time, most prior approaches use these signals to refine hidden representations, reweight attention, or fuse side information. That is, they acknowledge multi-scope behavior, but they do not usually turn scope differences into an explicit routing design problem. RouteRec builds on this literature by asking not only how to represent multiple scopes, but also how to assign computation differently across them. Its coarse-to-fine multi-stage routing therefore treats macro, mid, and micro evidence as distinct control views rather than as one pooled feature stream.

\subsection{Expert Routing Semantics and Inspectability}
MoE-based recommendation models show that expert specialization can improve flexibility, but routing semantics often remain unclear. In many architectures, routing is driven by hidden representations, task identities, or learned latent facets \cite{ma2018mmoe,tang2020ple,liu2025fame,bian2023m3srec,zhang2025hm4sr,zhou2022expertchoice,zoph2022stmoe}. This can be effective, yet it leaves a limited explanation of why a particular expert was chosen and whether that choice corresponds to an observable behavioral regime.

For our setting, this interpretability issue matters because RouteRec is not only trying to add capacity; it is trying to specialize computation according to behavior. We therefore separate the role of representations and routing. Hidden states remain the content channel transformed by experts, while routing is controlled by portable behavioral cues derived from standard logs. Prior expert models mainly specialize through latent, task-level, or facet-level signals, whereas RouteRec specializes through observable cues tied to temporal behavioral regimes. This behavior-guided routing view makes expert allocation easier to inspect, compare across stages, and analyze in later ablations.

\begin{figure*}[!t]
  \centering
  \includegraphics[width=0.95\textwidth]{figures/figure2_overall_arch_simple.pdf}
  \caption{Overall RouteRec framework. (a) RouteRec keeps a strong sequential backbone and replaces shared feed-forward blocks with behavior-guided MoE stages. (b) All stages use the same scalar-group conditional routing composition: a group-level distribution and a conditional intra-group distribution are multiplied to form expert probabilities. (c) Macro, mid, and micro differ by temporal evidence scope and routing granularity, not by router formula.}
  \Description{A three-panel architecture figure showing RouteRec with a unified scalar-group conditional router across macro, mid, and micro stages, where stages differ by temporal evidence scope and granularity.}
  \label{fig:framework}
\end{figure*}

\section{Method}
\label{sec:method}

\subsection{RouteRec Overview}
\label{sec:problem-framework}

Figure~\ref{fig:framework} summarizes RouteRec. For each user, we observe interactions ordered by time and train the model for next-item prediction. In one training instance, the observed prefix is denoted by $x_{1:t}$ and the target by $x_{t+1}$. Each interaction contains at least an item identity and timestamp, and may additionally include coarse category information depending on the dataset. RouteRec retains a strong sequential backbone for encoding the prefix, but replaces shared feed-forward computation with stage-wise MoE routing.

The model uses two signal channels. The representation channel maps the observed prefix to hidden states, which are the content representations transformed by experts. The behavior channel summarizes the same prefix into compact cues derived from ordinary logs. These cues are deliberately lightweight: they are designed as a control interface for routing rather than as another deep side-information branch. We first describe RouteRec as a design framework and then refer to the fixed paper instantiation only when concrete implementation choices matter. Appendix~\ref{app:method-details} collects the corresponding specification details.

This design directly follows the challenge structure introduced earlier. Challenge~1 motivates portable behavior cues derived from widely available logs. Challenge~2 motivates stage-wise routing over macro, mid, and micro temporal scopes. Challenge~3 motivates decoupling behavior-grounded routing from the expert transformation pathway so that expert allocation can be analyzed in terms of observable behavior. The remainder of this section formalizes how these three ideas are instantiated in one RouteRec architecture.

Let $\mathcal{T}=\{\mathrm{macro},\mathrm{mid},\mathrm{micro}\}$ be the stage set. Let $\mathbf{h}^{\mathrm{in},(s)}_t$ denote the hidden state entering stage $s$. At stage $s \in \mathcal{T}$, routing uses a stage-specific behavior feature $\mathbf{f}^{(s)}_t$, while experts consume the current stage input state:
\begin{align}
\boldsymbol{\pi}^{(s)}_t
&=
\mathrm{softmax}\!\left(\mathrm{Router}_s(\mathbf{f}^{(s)}_t)/\tau_s\right), \\
\tilde{\mathbf{h}}^{(s)}_t
&=
\sum_{e=1}^{E}\pi^{(s)}_{t,e}\,\mathrm{Expert}^{(s)}_e(\mathbf{h}^{\mathrm{in},(s)}_t).
\label{eq:stage-routing}
\end{align}
Eq.~\eqref{eq:stage-routing} is the common stage interface. In the serial RouteRec architecture, each stage consumes the hidden state produced by the previous attention or routing block. Thus, the framework keeps one shared MoE interface, while the paper instantiation specifies how these stage-local states are connected in sequence.

\subsection{Portable Behavioral Cues}
\label{sec:feature-construction}

RouteRec is designed for realistic data availability. It uses signals that are commonly present in interaction logs: item identity, timestamp, optional category or theme labels, and corpus-level popularity statistics derived from the training data. We intentionally avoid heavy side-information encoders in the routing branch. Instead, raw signals are converted into compact scalar summaries that are easy to compute, stable across datasets, and interpretable at analysis time.

For session $s_m$ and target position $t$, the stage windows are
\begin{align}
\mathcal{R}^{\mathrm{macro}}_t &= [s_{\max(1,m-H)},\dots,s_{m-1}], \\
\mathcal{R}^{\mathrm{mid}}_t &= [x_1,\dots,x_{t-1}], \\
\mathcal{R}^{\mathrm{micro}}_t &= [x_{\max(1,t-w)},\dots,x_{t-1}],
\end{align}
where $H=5$ and $w=5$ in our main setting. A shared summary operator $\mathcal{A}(\cdot)$ maps each window to a low-dimensional feature vector:
\[
\mathbf{f}^{(s)}_t=\mathcal{A}(\mathcal{R}^{(s)}_t),\qquad s\in\mathcal{T}.
\]

The summary dimensions are grouped into four semantic families: \textit{Tempo}, \textit{Focus}, \textit{Memory}, and \textit{Exposure}. Tempo captures coverage, pace, and interval dynamics. Focus captures concentration and switching behavior over category or theme proxies. Memory captures repetition, overlap, and recurrence structure. Exposure captures the level, spread, and drift of interacted popularity. In the fixed paper instantiation, this family template yields a compact 16-dimensional cue bank per stage. Categories are used when available; otherwise the corresponding focus cues are simply omitted through the missing-field handling documented in Appendix~\ref{app:method-details}. Representative examples are listed in Table~\ref{tab:appendix-feature-families}.

Importantly, the feature template is shared across stages, but each stage reads it through a different temporal lens. Macro features reflect persistent tendency and broad preference shape across recent sessions. Mid features track evolving intent inside the current session. Micro features emphasize immediate cues around the prediction target. This shared-template, stage-specific-view design keeps the feature branch compact while preserving temporal specificity.

This feature design is intentionally more structured than generic feature fusion. The goal is not to expose as many handcrafted statistics as possible, but to construct a small cue bank whose dimensions remain semantically legible and transportable across datasets. We therefore use stage-specific summaries with a consistent family template, rather than dataset-specific deep side branches. The appendix records representative cue definitions and the fixed paper setting, while more exhaustive column-level formulas can remain outside the main text because they are engineering details rather than the conceptual core of RouteRec.

\subsection{Coarse-to-Fine Multi-Stage Routing}

Although the router formula is shared, stages are not redundant. They differ in temporal evidence, routing granularity, and placement inside the backbone. In the fixed paper instantiation, the serial layout is
\[
\text{self-attention} \rightarrow \text{macro routing block} \rightarrow \text{mid routing block} \rightarrow \text{self-attention} \rightarrow \text{micro routing block}.
\]
Let $\mathbf{h}^{(0)}_t$ denote the hidden state after the first attention block, $\mathbf{h}^{(1)}_t$ after macro routing, $\mathbf{h}^{(2)}_t$ after mid routing, and $\mathbf{h}^{(3)}_t$ after the second attention block. Macro, mid, and micro therefore consume $\mathbf{h}^{(0)}_t$, $\mathbf{h}^{(1)}_t$, and $\mathbf{h}^{(3)}_t$, respectively. The first attention layer builds a contextualized prefix representation before any routing decision is made. Macro and mid then operate consecutively because both provide slower sequence-level control over closely related contextual states: macro reflects broader recent-history tendency, whereas mid refines the evolving intent of the current session. We do not insert another attention layer between them because that extra update increases depth without introducing a comparably distinct temporal view. By contrast, the second attention layer is placed immediately before micro so that the final local router sees a refreshed token interaction state before prediction. Appendix~\ref{app:stage-layout} reports the corresponding stage-order and layout ablations.

\paragraph{Macro and mid (sequence-wise routing).}
Macro and mid compute one routing distribution per sequence from their respective windows, then broadcast it over valid positions inside the sequence. Macro uses broader cross-session history; mid uses within-session prefix evidence. This setting emphasizes stable allocation for slower-changing context.

\paragraph{Micro (position-wise routing).}
Micro computes routing per token near the prediction target using short-horizon context. This enables fast adaptation to immediate transitions without changing the router form itself.

Formally, let $\bar{\mathbf{f}}_i^{\,(s)}$ denote the sequence-level cue summary used for sequence-wise routing at stage $s$. Then macro and mid use
\[
\begin{aligned}
\bar{\boldsymbol{\pi}}_i^{\,(s)}
&=
\mathrm{softmax}\!\left(\mathrm{Router}_s(\bar{\mathbf{f}}_i^{\,(s)})/\tau_s\right),
\qquad s\in\{\mathrm{macro},\mathrm{mid}\}, \\
\boldsymbol{\pi}^{(s)}_{i,t}
&=
\bar{\boldsymbol{\pi}}_i^{\,(s)}.
\end{aligned}
\]
Micro instead applies the same router form directly to position-specific cues,
\[
\boldsymbol{\pi}^{(\mathrm{micro})}_{i,t}
=
\mathrm{softmax}\!\left(\mathrm{Router}_{\mathrm{micro}}(\mathbf{f}^{(\mathrm{micro})}_{i,t})/\tau_{\mathrm{micro}}\right).
\]
Hence, RouteRec's coarse-to-fine behavior comes from stage evidence scope and update frequency, while the underlying router composition remains constant. The stages are therefore not repeated copies of the same router; they correspond to different temporal horizons at which behavioral evidence becomes informative for computation allocation.

\subsection{Behavior-Guided Expert Allocation}

The final RouteRec architecture uses one shared scalar-group conditional composition at every stage. Experts are partitioned into $G$ semantic groups with $C$ experts per group, so $E=GC$. In the paper instantiation, the group structure is organized around the four cue families, while capacity is expanded by multiple experts within each group. This alignment defines a semantic prior over expert organization rather than a hard hand assignment of individual cues to individual experts. For stage $s$ and position $t$, the router produces
\begin{align}
p^{(s)}_{t,g} &= \mathrm{softmax}_g\!\left(z^{(s)}_{t,g}\right), \\
r^{(s)}_{t,g,c} &= \mathrm{softmax}_c\!\left(u^{(s)}_{t,g,c}\right),
\end{align}
where $p^{(s)}_{t,g}$ is a group-level routing distribution and $r^{(s)}_{t,g,c}$ is a conditional within-group routing distribution. The final expert probability is their product:
\begin{equation}
\pi^{(s)}_{t,g,c}=p^{(s)}_{t,g}\,r^{(s)}_{t,g,c},
\label{eq:group-conditional-product}
\end{equation}
followed by flattening $(g,c)$ to expert index $e$. Eq.~\eqref{eq:group-conditional-product} is the same for macro, mid, and micro.

In the main configuration, all three stages use the same feature-driven group-conditional router. A group-level scalar router first determines which semantic control view is most relevant. A conditional intra-group router then distributes probability mass among the experts inside that group. In plain terms, RouteRec first selects the family-level control view and then allocates capacity within that view. This is the core reason we describe the method as behavior-guided allocation: the router is controlled by the cue bank, while hidden states remain the objects transformed by the experts. We use dense soft mixing in the paper because it yields stable diagnostics and cleaner stage-wise usage comparisons \cite{zoph2022stmoe,zhou2022expertchoice}; sparse top-$k$ routing is left outside the current scope.

This unified composition is important for two reasons. First, it gives all stages a shared structural prior, which reduces stage-specific router engineering and simplifies debugging. Second, it keeps diagnostics comparable across stages because differences in routing behavior can be attributed to evidence scope rather than formula changes.

In RouteRec, routing is driven by behavior cues, while hidden states remain dedicated to expert transformation in Eq.~\eqref{eq:stage-routing}. This preserves the routing-representation decoupling principle that underlies our approach.

\subsection{Training Objective}
\label{sec:training-objective}

RouteRec is trained with a standard next-item objective and two router-side regularizers in the main paper configuration: route consistency and z-loss. The first regularizer shapes routing semantics, and the second stabilizes routing numerics.

Let $\mathbf{z}$ be final logits and $y$ be the target item. The main loss is
\[
\mathcal{L}_{\mathrm{CE}}
=
-\log \frac{\exp z_y}{\sum_{j}\exp z_j}.
\]
For route consistency, we first compute a sample-level routing distribution at each stage by averaging position-level routing weights:
\[
\bar{\boldsymbol{\pi}}_i^{\,s}
=
\frac{1}{L_i}\sum_{t=1}^{L_i}\boldsymbol{\pi}_{i,t}^{\,s}
\]
where $L_i$ is the valid sequence length. We also form a stage-level cue summary
\[
\bar{\mathbf{f}}_i^{\,s}
=
\frac{1}{L_i}\sum_{t=1}^{L_i}\mathbf{f}_{i,t}^{\,s}.
\]
For each stage $s$, we construct a pair set $\mathcal{P}_s$ from the $k_{\mathrm{nn}}$ nearest non-self samples in the mini-batch according to cosine similarity between these stage-level cue summaries, and minimize
\[
\mathcal{L}_{\mathrm{cons}}
=
\frac{1}{|\mathcal{T}|}
\sum_{s \in \mathcal{T}}
\frac{1}{|\mathcal{P}_s|}
\sum_{(i,j)\in \mathcal{P}_s}
\mathrm{JSD}\!\left(\bar{\boldsymbol{\pi}}_i^{\,s}, \bar{\boldsymbol{\pi}}_j^{\,s}\right),
\]
which encourages neighboring behavioral states to produce compatible routing distributions. In other words, semantically similar inputs should not induce arbitrarily different routing behavior. Appendix~\ref{app:method-details} reports the neighbor setting used in the paper instantiation.
For macro and mid, the sample-level average reduces to the broadcast sequence-wise route; we keep the averaging notation for a unified treatment across stages.

\begin{table*}[t]
  \centering
  \small
  \setlength{\tabcolsep}{3.0pt}
  \renewcommand{\arraystretch}{1.02}
  \resizebox{\textwidth}{!}{%
  \begin{tabular}{llrrrrrrrrrrr}
    \toprule
    Dataset & Metric & SASRec & GRU4Rec & TiSASRec & FEARec & DuoRec & BSARec & FAME & DIF-SR & FDSA & RouteRec & Improv. \\
    \midrule
    \multirow{3}{*}{Beauty}
      & HR@10   & \tblsecond{0.1347} & 0.0716 & 0.1318 & 0.1261 & 0.1146 & 0.0516 & 0.0315 & 0.1060 & \tblsecond{0.1347} & \tblbest{0.1616} & +20.0\% \\
      & NDCG@10 & 0.0748 & 0.0340 & 0.0878 & 0.0696 & 0.0680 & 0.0316 & 0.0211 & 0.0685 & \tblsecond{0.0909} & \tblbest{0.0982} & +8.1\% \\
      & MRR@20  & 0.0588 & 0.0247 & 0.0770 & 0.0544 & 0.0569 & 0.0260 & 0.0190 & 0.0592 & \tblsecond{0.0805} & \tblbest{0.0822} & +2.0\% \\
    \midrule
    \multirow{3}{*}{Retail Rocket}
      & HR@10   & \tblsecond{0.5605} & 0.5211 & 0.5435 & 0.5441 & 0.5377 & 0.4703 & 0.4697 & 0.4365 & 0.5133 & \tblbest{0.5612} & +0.1\% \\
      & NDCG@10 & \tblsecond{0.4099} & 0.3858 & 0.3963 & 0.3935 & 0.3864 & 0.3963 & 0.3934 & 0.3900 & 0.3651 & \tblbest{0.4150} & +1.2\% \\
      & MRR@20  & 0.3673 & 0.3479 & 0.3547 & 0.3515 & 0.3441 & \tblsecond{0.3743} & 0.3710 & \tblbest{0.3744} & 0.3239 & 0.3737 & -0.2\% \\
    \midrule
    \multirow{3}{*}{Foursquare}
      & HR@10   & \tblbest{0.3233} & 0.2130 & 0.3157 & 0.3042 & 0.3197 & 0.2462 & 0.2546 & 0.3034 & 0.2998 & \tblsecond{0.3225} & -0.2\% \\
      & NDCG@10 & 0.2014 & 0.1455 & \tblsecond{0.2016} & 0.1925 & 0.1955 & 0.1639 & 0.1646 & 0.1984 & 0.2001 & \tblbest{0.2041} & +1.3\% \\
      & MRR@20  & 0.1679 & 0.1278 & 0.1700 & 0.1619 & 0.1607 & 0.1410 & 0.1401 & 0.1690 & \tblbest{0.1723} & \tblsecond{0.1712} & -0.6\% \\
    \midrule
    \multirow{3}{*}{ML-1M}
      & HR@10   & 0.1478 & 0.1413 & 0.1538 & 0.1487 & 0.1520 & \tblsecond{0.1589} & 0.1557 & 0.1501 & \tblbest{0.1594} & 0.1561 & -2.0\% \\
      & NDCG@10 & 0.0770 & 0.0742 & 0.0844 & 0.0724 & 0.0811 & \tblbest{0.0854} & \tblsecond{0.0853} & 0.0764 & 0.0835 & 0.0771 & -9.7\% \\
      & MRR@20  & 0.0608 & 0.0587 & \tblbest{0.0690} & 0.0550 & 0.0649 & 0.0678 & \tblsecond{0.0687} & 0.0595 & 0.0658 & 0.0585 & -15.3\% \\
    \midrule
    \multirow{3}{*}{LastFM$^{\dagger}$}
      & HR@10   & \tblbest{0.3994} & 0.3090 & 0.3780 & 0.3765 & \tblsecond{0.3831} & 0.3453 & 0.3312 & 0.3617 & 0.3657 & 0.3794 & -5.0\% \\
      & NDCG@10 & \tblbest{0.3267} & 0.2673 & 0.3196 & 0.3167 & 0.3083 & 0.3024 & 0.2900 & 0.3149 & 0.3171 & \tblsecond{0.3253} & -0.5\% \\
      & MRR@20  & \tblsecond{0.3054} & 0.2561 & 0.3027 & 0.2989 & 0.2863 & 0.2902 & 0.2784 & 0.3020 & 0.3033 & \tblbest{0.3099} & +1.5\% \\
    \midrule
    \multirow{3}{*}{KuaiRec$^{\dagger}$}
      & HR@10   & \tblsecond{0.1055} & 0.0945 & 0.1011 & 0.1005 & 0.0995 & 0.0918 & 0.1011 & 0.0918 & 0.0967 & \tblbest{0.1060} & +0.5\% \\
      & NDCG@10 & \tblsecond{0.0943} & 0.0790 & 0.0916 & 0.0919 & 0.0919 & 0.0890 & 0.0915 & 0.0871 & 0.0865 & \tblbest{0.0965} & +2.3\% \\
      & MRR@20  & \tblsecond{0.0922} & 0.0753 & 0.0894 & 0.0903 & 0.0905 & 0.0882 & 0.0899 & 0.0863 & 0.0843 & \tblbest{0.0948} & +2.7\% \\
    \bottomrule
  \end{tabular}}
      \caption{Main experimental results. For each baseline, we report the test metrics of the run achieving the highest mean validation score over the nine reported metrics. RouteRec is reported from the best candidate run under the same selection criterion. $^{\dagger}$ Reported on the processed subsets used in our runs: 20\% for KuaiRec and 3\% for LastFM. Bold and underlined values denote the best and one second-best result within each row.}
  \label{tab:main-overall}
\end{table*}
To stabilize routing numerics, we adopt z-loss on the pre-softmax expert logits $a_{i,t,e}^{\,s}$~\cite{zoph2022stmoe}:
\[
\mathcal{L}_{z}
=
\frac{1}{|\mathcal{T}|}
\sum_{s \in \mathcal{T}}
\frac{1}{N_s}
\sum_{(i,t)\in\mathcal{V}_s}
\left(
\log \sum_{e=1}^{E}\exp a_{i,t,e}^{\,s}
\right)^2,
\]
where $\mathcal{V}_s$ is the set of valid positions and $N_s=|\mathcal{V}_s|$. This term discourages excessive logit scale and helps prevent unstable, overly sharp routing. The final objective is
\[
\mathcal{L}
=
\mathcal{L}_{\mathrm{CE}}
+
\lambda_{\mathrm{cons}}\,\mathcal{L}_{\mathrm{cons}}
+
\lambda_{z}\,\mathcal{L}_{z},
\]
where $\lambda_{\mathrm{cons}}$ and $\lambda_z$ are trade-off coefficients. Unlike classical MoE setups that explicitly enforce uniform expert traffic, the selected main configuration does not include a load-balancing term in the final paper runs. Our goal here is not uniform usage for its own sake, but stable and behavior-aligned specialization. Appendix~\ref{app:objective} reports objective variants including optional balance regularization.

\section{Experiments}
\label{sec:experiments}


\subsection{Experimental Setup}
We evaluate RouteRec on six public datasets. These datasets cover e-commerce, location recommendation, short-video consumption, music listening, movie preference, and anonymous clickstream behavior. The six corpora are Beauty, Foursquare, KuaiRec, LastFM, ML-1M, and Retail Rocket. Appendix~\ref{app:data-details} summarizes their provenance and processed statistics.

Across all datasets, interactions are first ordered chronologically and then partitioned into sessions using dataset-specific inactivity thresholds. We adopt a 30-minute threshold for Foursquare, KuaiRec, LastFM, ML-1M, and Retail Rocket, whereas Beauty uses a 14-day threshold because review events are substantially sparser than click, play, or check-in logs. After sessionization, we perform iterative $5$-core-style filtering until convergence: sessions of length below $5$ are removed, items with fewer than $3$ interactions are discarded, and session lengths are re-evaluated after item pruning. Beauty is constructed from Amazon Beauty reviews and additionally excludes interactions with rating below $3$.

For transparency, the reported KuaiRec and LastFM results are obtained from processed subsets rather than the full raw corpora: specifically, we use a 20\% KuaiRec subset and a 3\% LastFM subset. Table~\ref{tab:appendix-datasets} therefore reports statistics for the final processed datasets used in the experiments. After preprocessing, sessions are sorted by start time and divided at the session level into 70\%/15\%/15\% train/validation/test partitions. In the initial split, the earliest 70\% of sessions are assigned to training, and the remaining tail is divided into validation and test using stratification over target-repeat status, target-popularity bucket, and session-length bucket. For the final evaluation files, the training split is kept fixed, the held-out sessions are merged and re-ordered by session start time, and the merged tail is then divided contiguously into equal validation and test halves. During this held-out cleanup step, rows whose items are unseen in training are removed unless the unseen item is the final target of the session; unseen target cases are retained. When category or time signals are incomplete, RouteRec falls back to the same portable Tempo, Focus, Memory, and Exposure feature families computed from the available interaction log.

The baselines include GRU4Rec, SASRec, TiSASRec, DuoRec, FEARec, BSARec, FAME, DIF-SR, and FDSA~\cite{hidasi2016gru4rec,kang2018sasrec,li2020tisasrec,qiu2022duorec,du2023fearec,shin2024bsarec,liu2025fame,xie2022difsr,zhang2019fdsa}. All methods are trained within the same optimization family, with matched early-stopping rules and comparable capacity ranges. In the comparison tables, each baseline is represented by the run attaining the highest mean validation score over the nine reported metrics, and we report the corresponding test metrics. RouteRec is selected from the candidate runs considered under the same criterion. This reporting protocol follows recent calls for transparent and controlled comparison in session-based recommendation \cite{ludewig2021empirical}.

\subsection{Overall Performance: Does Behavior-Guided Routing Help?}
\label{sec:overall-performance}
Table~\ref{tab:main-overall} presents the main experimental comparison in the text using HR@10, NDCG@10, and MRR@20. This presentation keeps the main paper focused on a compact set of representative cutoffs, while Appendix~\ref{app:data-details} reports the full $k\in\{5,10,20\}$ grid. RouteRec exhibits its clearest gains on Beauty and the KuaiRec split, remains competitive on Retail Rocket, and stays close to the strongest shared-backbone baselines on Foursquare, LastFM, and ML-1M. This pattern is consistent with the dataset characterization in Appendix~\ref{app:data-details}: datasets with stronger short-horizon or concentrated behavioral signals provide clearer scope-separated routing cues than more stable preference-driven corpora.

Two points are worth emphasizing. First, the improvement pattern is strongest when the target slice is behaviorally concentrated, which matches the motivation for behavior-guided routing. Second, the gain does not reduce to simple capacity scaling: several baselines already use rich sequence encoders or feature channels, yet RouteRec still benefits when the conditional computation path is chosen from explicit behavior cues rather than shared uniformly.

\subsection{What Should Control Routing?}
\label{sec:routing-control}
Figure~\ref{fig:routing-control-panels} compares four routing controls: a shared FFN, hidden-only routing, mixed hidden+behavior routing, and the final behavior-guided router. Hidden-only routing can improve over a shared FFN, but the gain is smaller and less stable than behavior-guided routing. Re-injecting hidden states into the router also fails to produce a consistent advantage over the behavior-only design.

The consistency panel makes the structural reason clearer. When routing is tied to observable cues rather than hidden-state geometry alone, behaviorally similar cases are assigned more compatible expert patterns. The advantage of RouteRec therefore comes from grounding the routing decision in semantically readable evidence, not from adding a more complicated gate.

\begin{figure}[!t]
  \centering
  \begin{minipage}[t]{0.48\columnwidth}
    \centering
    \placeholderpanel{1.25in}

    \vspace{0.35em}
    {\footnotesize (a) Avg. ranking quality across routing controls}
  \end{minipage}\hfill
  \begin{minipage}[t]{0.48\columnwidth}
    \centering
    \placeholderpanel{1.25in}

    \vspace{0.35em}
    {\footnotesize (b) Route consistency over feature-similarity buckets}
  \end{minipage}
  \caption{Routing-control comparison. (a) Average MRR@20 across shared FFN, hidden-only, mixed, and behavior-guided routing. (b) Routing similarity over feature-similarity buckets.}
  \Description{A two-panel placeholder figure comparing routing-control variants by average ranking quality and routing consistency over feature-similarity buckets.}
  \label{fig:routing-control-panels}
\end{figure}

\subsection{Why Is the Staged RouteRec Structure Effective?}
\label{sec:stage-structure}
Figure~\ref{fig:stage-structure-panels} examines the RouteRec architecture itself. Panel (a) removes macro, mid, or micro routing to isolate the role of each temporal scope. Panel (b) separates staged routing from generic FFN replacement by comparing dense replacement against one-, two-, and three-stage layouts. Panel (c) tests router-composition and stage-order variants to show that the final composition is not an arbitrary design choice.

Across all three views, the same qualitative pattern appears. Performance improves when routing is distributed across temporal scopes, dense replacement alone does not recover the same gain, and the final router composition remains among the strongest variants. The benefit of RouteRec therefore comes from coarse-to-fine behavioral allocation rather than from expert insertion by itself.

\begin{figure*}[!t]
  \centering
  \begin{minipage}[t]{0.31\textwidth}
    \centering
    \placeholderpanel{1.05in}

    \vspace{0.35em}
    {\footnotesize (a) Stage removal}
  \end{minipage}\hfill
  \begin{minipage}[t]{0.31\textwidth}
    \centering
    \placeholderpanel{1.05in}

    \vspace{0.35em}
    {\footnotesize (b) Dense replacement vs. staged routing}
  \end{minipage}\hfill
  \begin{minipage}[t]{0.31\textwidth}
    \centering
    \placeholderpanel{1.05in}

    \vspace{0.35em}
    {\footnotesize (c) Router-composition and order variants}
  \end{minipage}
  \caption{Stage and layout ablations. (a) Stage removal. (b) Dense replacement versus staged routing. (c) Router-composition and stage-order variants.}
  \Description{A three-panel placeholder figure for stage removal, dense-versus-staged routing comparisons, and router-composition or stage-order variants.}
  \label{fig:stage-structure-panels}
\end{figure*}

\subsection{Are Lightweight Cues Enough?}
\label{sec:lightweight-cues}
Figure~\ref{fig:cue-ablation-panels} evaluates whether RouteRec depends on rich metadata or whether ordinary interaction logs already provide enough routing evidence. The left panel removes cue families from the full model, and the right panel normalizes each variant by the gain of the full-cue model. Removing category-derived or time-derived cues lowers performance, but the sequence-only variant still retains a substantial fraction of the improvement over the shared-backbone baseline.

This result strengthens the portability claim. RouteRec benefits from richer evidence when it is available, but it does not require a heavy side-information stack. A meaningful part of the routing gain can already be recovered from the signals that most sequential recommendation systems already log.

\begin{figure}[!t]
  \centering
  \begin{minipage}[t]{0.48\columnwidth}
    \centering
    \placeholderpanel{1.25in}

    \vspace{0.35em}
    {\footnotesize (a) Avg. MRR@20 by cue family removal}
  \end{minipage}\hfill
  \begin{minipage}[t]{0.48\columnwidth}
    \centering
    \placeholderpanel{1.25in}

    \vspace{0.35em}
    {\footnotesize (b) Relative retention of the full-cue gain}
  \end{minipage}
  \caption{Cue-family ablations. (a) Average MRR@20 under cue removal. (b) Gain retention relative to the full cue set.}
  \Description{A two-panel placeholder figure for cue-family ablations and normalized gain retention.}
  \label{fig:cue-ablation-panels}
\end{figure}

Detailed method notes, full-cutoff results, and extended analyses, including slice-wise and transfer studies, are reported in Appendices~\ref{app:method-details}--\ref{app:transfer}.

\section{Conclusion}
\label{sec:conclusion}
RouteRec starts from a simple hypothesis: when behavioral evidence changes, the computation path should change as well. The experiments support this view from three main angles. First, RouteRec improves next-item ranking over strong sequential baselines across most datasets and cutoffs. Second, the gain comes from the routing signal and the staged design rather than from inserting experts indiscriminately. Third, lightweight cues derived from ordinary interaction logs already recover a substantial share of the benefit, which keeps the method practical even when side information is limited.

Taken together, these results suggest that compact behavioral cues provide a useful control interface for conditional computation in recommendation. Future work can build on this interface with broader transfer regimes, richer but still robust cue construction, and tighter deployment-aware constraints such as latency, continual adaptation, and router calibration.

\bibliographystyle{ACM-Reference-Format}
\bibliography{refs}

\nobalance
\clearpage
\appendix

\section{Method Details and Design Rationale}
\label{app:method-details}
This appendix records the fixed paper instantiation and the representative cue bank used by RouteRec. The main text emphasizes the conceptual design axes, whereas this section captures the implementation choices that are important for reproducibility and for interpreting later ablations.

\subsection{Main Paper Configuration}

\begin{table}[t]
  \centering
  \caption{Fixed main configuration used in the paper.}
  \label{tab:appendix-main-setting}
  \footnotesize
  \setlength{\tabcolsep}{4pt}
  \begin{tabularx}{\columnwidth}{@{}>{\raggedright\arraybackslash}X>{\raggedright\arraybackslash}X@{}}
    \toprule
    Aspect & Main paper configuration \\
    \midrule
    Backbone layout & self-attention $\rightarrow$ macro routing block $\rightarrow$ mid routing block $\rightarrow$ self-attention $\rightarrow$ micro routing block \\
    Stage execution & Serial \\
    Routing granularity & Macro and mid use sequence-wise routing; micro uses position-wise routing \\
    Router composition & Shared scalar-group conditional router at all three stages \\
    Expert mixing & Dense softmax mixture over the full expert set \\
    Macro history window & $H=5$ recent sessions \\
    Core auxiliary terms & Route consistency + z-loss; load balancing disabled in the selected paper runs \\
    \bottomrule
  \end{tabularx}
\end{table}

The placement of the two attention layers reflects the intended temporal roles of the stages. The first attention layer contextualizes the raw prefix before any routing is applied. Macro and mid then operate consecutively because both provide slower sequence-level control over the same contextualized sequence state: macro captures broader recent-history tendency, while mid refines within-session evolution. A second attention layer is inserted only before micro so that the final local router receives a refreshed token interaction state immediately before prediction. This layout was the most stable high-performing architecture in our architecture sweep, and Appendix~\ref{app:stage-layout} compares it against denser and differently ordered alternatives.

The table below should be read as a reproducibility aid rather than as a search-space summary. Architectural choices are fixed across the paper track, whereas capacity and optimization values are selected per dataset from a narrow candidate set and are therefore reported as compact ranges only when no single global value is shared across all reported runs.

\subsection{Implementation Specification}

\begin{table}[t]
  \centering
  \caption{Implementation details reported for reproducibility.}
  \label{tab:appendix-implementation-spec}
  \footnotesize
  \setlength{\tabcolsep}{4pt}
  \begin{tabularx}{\columnwidth}{@{}>{\raggedright\arraybackslash}X>{\raggedright\arraybackslash}X@{}}
    \toprule
    Item & Reported setting \\
    \midrule
    Semantic groups & $G=4$ groups aligned with Tempo, Focus, Memory, and Exposure \\
    Experts per group & $C\in\{2,3,4\}$ in the paper track; total experts $E=4C$ \\
    Cue dimensionality & $16$ scalars per stage in the fixed paper instantiation \\
    Router hidden width & Selected from a small candidate set in the paper track; representative range $48$--$128$ \\
    Routing temperature & Macro $1.0$; mid and micro $1.2$ in the reported setting \\
    Consistency neighbors & $k_{\mathrm{nn}}=1$ nearest non-self sample per stage in the paper track \\
    Missing-signal handling & Broadcast scalar fields across positions when available; otherwise use zero filling for the missing sequence field \\
    \bottomrule
  \end{tabularx}
\end{table}

\subsection{Representative Behavioral Cue Families}

\begin{table*}[t]
  \centering
  \caption{Representative cue summaries used by RouteRec in the main configuration. Each stage uses four semantic families with four scalar summaries per family; the table lists representative examples rather than every implementation detail.}
  \label{tab:appendix-feature-families}
  \footnotesize
  \setlength{\tabcolsep}{4pt}
  \renewcommand{\arraystretch}{1.05}
  \begin{tabularx}{\textwidth}{@{}lXXX@{}}
    \toprule
    Family & Macro (recent-session tendency) & Mid (within-session evolution) & Micro (short-horizon pressure) \\
    \midrule
    Tempo & Context-valid ratio, last inter-session gap, mean pace, pace trend & Valid-prefix ratio, mean/std interaction interval, session age & Last gap, short-window mean gap, short-window std, deviation from mid-scale pace \\
    Focus & Theme entropy, dominant-theme mass, theme repeat rate, theme shift rate & Category entropy, category top-1 mass, category switch rate, category uniqueness & Immediate category switch, mismatch to recent category history, suffix entropy, suffix uniqueness \\
    Memory & Repeat intensity, adjacent category overlap, adjacent item overlap, repeat trend & Item uniqueness, repeat rate, novelty rate, longest item run & Reconsumption flag, suffix reconsumption rate, suffix item uniqueness, suffix longest run \\
    Exposure & Mean popularity, popularity spread, popularity entropy, popularity trend & Mean popularity, popularity std, popularity entropy, popularity drift & Last-item popularity, suffix popularity spread, suffix popularity entropy, deviation from mid-scale popularity \\
    \bottomrule
  \end{tabularx}
\end{table*}

This cue bank is intentionally compact. The feature families are hand-designed in the sense that they encode behavioral hypotheses, but they are not meant to become an unbounded feature-engineering exercise. Instead, the design goal is to build a small set of summaries that can act as a stable routing interface across datasets. In practice, the macro bank is instantiated from the recent-session window used by the selected paper instantiation, while mid and micro are computed from the current session prefix at different temporal scopes. The main text therefore reports the family structure and representative examples, while this appendix carries the more implementation-oriented view.

\section{Dataset, Evaluation, and Full Results}
\label{app:data-details}
This appendix summarizes the six datasets and the common evaluation protocol. The original releases are cited here for completeness~\cite{mcauley2015amazon,yang2010foursquare,gao2022kuairec,celma2010music,harper2015movielens,retailrocket2015dataset}. The dataset descriptions emphasize domain differences and signal availability because the routing branch depends on which behavioral evidence is exposed in each corpus.

\subsection{Datasets and Protocol}

\begin{table}[t]
  \centering
  \footnotesize
  \setlength{\tabcolsep}{3.0pt}
  \begin{tabular*}{\columnwidth}{@{\extracolsep{\fill}}llrrrc@{}}
    \toprule
    Dataset & Domain & Interactions & Sessions & Items & Avg. len. \\
    \midrule
    Beauty & E-com. & 33,488 & 4,243 & 3,625 & 7.9 \\
    Foursquare & POI & 145,238 & 25,369 & 30,588 & 5.7 \\
    KuaiRec$^{\dagger}$ & Video & 287,411 & 24,458 & 6,477 & 11.8 \\
    LastFM$^{\dagger}$ & Music & 470,408 & 25,089 & 52,510 & 18.8 \\
    ML-1M & Movie & 575,281 & 14,539 & 3,533 & 39.6 \\
    Retail Rocket & Browse & 821,243 & 153,092 & 90,211 & 5.4 \\
    \bottomrule
  \end{tabular*}
  \caption{Datasets used in the experiments. $^{\dagger}$ Reported on the processed subsets used in our runs: 20\% for KuaiRec and 3\% for LastFM.}
  \label{tab:appendix-datasets}
\end{table}

Table~\ref{tab:appendix-datasets} reports the statistics of the final processed datasets used in the experiments rather than the raw corpus sizes. In particular, KuaiRec corresponds to the 20\% processed subset used in our runs, and LastFM corresponds to the 3\% processed subset.

Beauty and Retail Rocket stress sparse e-commerce revisit behavior, Foursquare and KuaiRec emphasize stronger short-horizon intent changes, and LastFM and ML-1M provide longer preference-driven histories. This spread is useful because it tests whether behavior-guided routing helps only in highly dynamic settings or remains competitive when preference structure is more stable.

Signal availability also differs in a way that matters for routing. Beauty exposes item, time, and product-category cues after low-rating removal. Foursquare includes item, timestamp, and venue-category structure with bursty check-in dynamics. KuaiRec offers item, time, and coarse content-side metadata in short-video sessions with strong short-term intent shifts. LastFM provides repeated listening histories with artist-side signals, while ML-1M contributes longer preference-driven sequences with genre information. Retail Rocket is the sparsest side-information setting, relying mainly on anonymous event streams with item and time signals.

The preprocessing pipeline is unified across datasets in overall structure, while the session boundaries are adapted to the logging regime of each corpus. Foursquare, KuaiRec, LastFM, ML-1M, and Retail Rocket are sessionized with a 30-minute inactivity threshold, whereas Beauty uses a 14-day threshold because review timestamps are substantially more diffuse. After sessionization, we repeatedly remove sessions shorter than $5$ and items with fewer than $3$ interactions until no further pruning is required. This iterative filtering is necessary because item removal can render previously valid sessions too short, making a single-pass procedure insufficient.

Session splits are constructed at the session level after sorting by session start time. The earliest 70\% of sessions are assigned to training, and the remaining 30\% are treated as a held-out tail. This tail is first divided into validation and test using stratification over three factors: whether the target item already appears in the prefix, the target item's popularity bucket with respect to the training split, and the session-length bucket. We then build the final evaluation files by merging the held-out sessions, re-ordering them by session start time, and re-splitting them contiguously into equal validation and test halves. In these held-out files, non-target rows whose items are unseen in training are removed, whereas unseen target items are preserved. We report HR, NDCG, and MRR at $k\in\{5,10,20\}$ for all methods.

For the result tables, each baseline is represented by the run with the highest mean validation score over the nine reported metrics, and RouteRec is reported from the best run among the candidate runs under the same selection rule.

\subsection{Full Ranking Results}
Table~\ref{tab:appendix-full-seen} restores the full cutoff grid summarized more compactly in Table~\ref{tab:main-overall}. This presentation preserves the same selection rule as the main text while exposing the complete $k\in\{5,10,20\}$ breakdown.

\begin{table*}[t]
  \centering
  \scriptsize
  \setlength{\tabcolsep}{2.6pt}
  \renewcommand{\arraystretch}{0.96}
  \resizebox{\textwidth}{!}{%
  \begin{tabular}{llrrrrrrrrrrr}
    \toprule
    Dataset & Metric & SASRec & GRU4Rec & TiSASRec & FEARec & DuoRec & BSARec & FAME & DIF-SR & FDSA & RouteRec & Improv. \\
    \midrule
    \multirow{9}{*}{Beauty}
      & HR@5    & \tblsecond{0.0974} & 0.0372 & 0.0974 & 0.0974 & 0.0974 & 0.0430 & 0.0258 & 0.0802 & 0.0946 & \tblbest{0.1054} & +8.2\% \\
      & HR@10   & \tblsecond{0.1347} & 0.0716 & 0.1318 & 0.1261 & 0.1146 & 0.0516 & 0.0315 & 0.1060 & \tblsecond{0.1347} & \tblbest{0.1616} & +20.0\% \\
      & HR@20   & \tblsecond{0.1719} & 0.1032 & 0.1719 & 0.1719 & 0.1719 & 0.0602 & 0.0516 & 0.1433 & 0.1719 & \tblbest{0.2014} & +17.2\% \\
      & NDCG@5  & 0.0629 & 0.0227 & 0.0765 & 0.0603 & 0.0625 & 0.0289 & 0.0192 & 0.0602 & \tblsecond{0.0779} & \tblbest{0.0803} & +3.1\% \\
      & NDCG@10 & 0.0748 & 0.0340 & 0.0878 & 0.0696 & 0.0680 & 0.0316 & 0.0211 & 0.0685 & \tblsecond{0.0909} & \tblbest{0.0982} & +8.1\% \\
      & NDCG@20 & 0.0842 & 0.0418 & 0.0980 & 0.0810 & 0.0827 & 0.0339 & 0.0260 & 0.0778 & \tblsecond{0.1003} & \tblbest{0.1083} & +8.0\% \\
      & MRR@5   & 0.0514 & 0.0179 & 0.0694 & 0.0476 & 0.0506 & 0.0242 & 0.0170 & 0.0533 & \tblbest{0.0725} & \tblsecond{0.0721} & -0.6\% \\
      & MRR@10  & 0.0563 & 0.0226 & 0.0742 & 0.0515 & 0.0528 & 0.0253 & 0.0178 & 0.0567 & \tblsecond{0.0779} & \tblbest{0.0794} & +1.9\% \\
      & MRR@20  & 0.0588 & 0.0247 & 0.0770 & 0.0544 & 0.0569 & 0.0260 & 0.0190 & 0.0592 & \tblsecond{0.0805} & \tblbest{0.0822} & +2.0\% \\
    \midrule
    \multirow{9}{*}{Retail Rocket}
      & HR@5    & \tblsecond{0.4869} & 0.4475 & 0.4700 & 0.4647 & 0.4611 & 0.4413 & 0.4365 & 0.4305 & 0.4344 & \tblbest{0.4883} & +0.3\% \\
      & HR@10   & \tblsecond{0.5605} & 0.5211 & 0.5435 & 0.5441 & 0.5377 & 0.4703 & 0.4697 & 0.4365 & 0.5133 & \tblbest{0.5612} & +0.1\% \\
      & HR@20   & \tblbest{0.6316} & 0.5863 & 0.6135 & 0.6200 & 0.6134 & 0.4986 & 0.5003 & 0.4409 & 0.5886 & \tblsecond{0.6313} & -0.1\% \\
      & NDCG@5  & 0.3860 & 0.3620 & 0.3724 & 0.3678 & 0.3616 & 0.3869 & 0.3826 & \tblsecond{0.3881} & 0.3396 & \tblbest{0.3914} & +0.9\% \\
      & NDCG@10 & \tblsecond{0.4099} & 0.3858 & 0.3963 & 0.3935 & 0.3864 & 0.3963 & 0.3934 & 0.3900 & 0.3651 & \tblbest{0.4150} & +1.2\% \\
      & NDCG@20 & \tblsecond{0.4280} & 0.4023 & 0.4139 & 0.4127 & 0.4057 & 0.4034 & 0.4012 & 0.3911 & 0.3842 & \tblbest{0.4327} & +1.1\% \\
      & MRR@5   & 0.3524 & 0.3335 & 0.3399 & 0.3355 & 0.3284 & \tblsecond{0.3685} & 0.3644 & \tblbest{0.3733} & 0.3082 & 0.3591 & -3.8\% \\
      & MRR@10  & 0.3623 & 0.3434 & 0.3498 & 0.3462 & 0.3388 & \tblsecond{0.3724} & 0.3689 & \tblbest{0.3741} & 0.3187 & 0.3688 & -1.4\% \\
      & MRR@20  & 0.3673 & 0.3479 & 0.3547 & 0.3515 & 0.3441 & \tblsecond{0.3743} & 0.3710 & \tblbest{0.3744} & 0.3239 & 0.3737 & -0.2\% \\
    \midrule
    \multirow{9}{*}{Foursquare}
      & HR@5    & \tblsecond{0.2518} & 0.1743 & 0.2474 & 0.2418 & 0.2474 & 0.2014 & 0.2014 & 0.2426 & 0.2478 & \tblbest{0.2562} & +1.7\% \\
      & HR@10   & \tblbest{0.3233} & 0.2130 & 0.3157 & 0.3042 & 0.3197 & 0.2462 & 0.2546 & 0.3034 & 0.2998 & \tblsecond{0.3225} & -0.2\% \\
      & HR@20   & \tblbest{0.3833} & 0.2590 & 0.3721 & 0.3645 & 0.3717 & 0.2866 & 0.3034 & 0.3521 & 0.3477 & \tblsecond{0.3777} & -1.5\% \\
      & NDCG@5  & 0.1782 & 0.1331 & 0.1793 & 0.1721 & 0.1721 & 0.1496 & 0.1475 & 0.1786 & \tblbest{0.1834} & \tblsecond{0.1827} & -0.4\% \\
      & NDCG@10 & 0.2014 & 0.1455 & \tblsecond{0.2016} & 0.1925 & 0.1955 & 0.1639 & 0.1646 & 0.1984 & 0.2001 & \tblbest{0.2041} & +1.3\% \\
      & NDCG@20 & \tblsecond{0.2167} & 0.1572 & 0.2159 & 0.2077 & 0.2088 & 0.1741 & 0.1769 & 0.2107 & 0.2123 & \tblbest{0.2181} & +0.7\% \\
      & MRR@5   & 0.1540 & 0.1195 & 0.1568 & 0.1492 & 0.1473 & 0.1324 & 0.1297 & 0.1574 & \tblbest{0.1620} & \tblsecond{0.1584} & -2.2\% \\
      & MRR@10  & 0.1636 & 0.1246 & 0.1660 & 0.1578 & 0.1569 & 0.1382 & 0.1368 & 0.1656 & \tblbest{0.1689} & \tblsecond{0.1673} & -0.9\% \\
      & MRR@20  & 0.1679 & 0.1278 & 0.1700 & 0.1619 & 0.1607 & 0.1410 & 0.1401 & 0.1690 & \tblbest{0.1723} & \tblsecond{0.1712} & -0.6\% \\
    \midrule
    \multirow{9}{*}{ML-1M}
      & HR@5    & 0.0957 & 0.0906 & 0.1018 & 0.0883 & 0.0990 & \tblbest{0.1050} & 0.0976 & 0.0888 & \tblsecond{0.1046} & 0.0939 & -10.6\% \\
      & HR@10   & 0.1478 & 0.1413 & 0.1538 & 0.1487 & 0.1520 & \tblsecond{0.1589} & 0.1557 & 0.1501 & \tblbest{0.1594} & 0.1561 & -2.0\% \\
      & HR@20   & 0.2254 & 0.2119 & \tblsecond{0.2375} & 0.2277 & 0.2305 & 0.2282 & 0.2212 & 0.2300 & \tblbest{0.2430} & 0.2337 & -3.8\% \\
      & NDCG@5  & 0.0605 & 0.0581 & \tblsecond{0.0677} & 0.0531 & 0.0640 & \tblbest{0.0680} & 0.0667 & 0.0566 & 0.0656 & 0.0572 & -15.9\% \\
      & NDCG@10 & 0.0770 & 0.0742 & 0.0844 & 0.0724 & 0.0811 & \tblbest{0.0854} & \tblsecond{0.0853} & 0.0764 & 0.0835 & 0.0771 & -9.7\% \\
      & NDCG@20 & 0.0965 & 0.0919 & \tblbest{0.1054} & 0.0924 & 0.1009 & 0.1028 & 0.1018 & 0.0963 & \tblsecond{0.1043} & 0.0964 & -8.5\% \\
      & MRR@5   & 0.0489 & 0.0474 & \tblsecond{0.0565} & 0.0416 & 0.0525 & 0.0559 & \tblbest{0.0566} & 0.0460 & 0.0529 & 0.0452 & -20.1\% \\
      & MRR@10  & 0.0556 & 0.0539 & \tblsecond{0.0633} & 0.0495 & 0.0595 & 0.0631 & \tblbest{0.0642} & 0.0542 & 0.0603 & 0.0534 & -16.8\% \\
      & MRR@20  & 0.0608 & 0.0587 & \tblbest{0.0690} & 0.0550 & 0.0649 & 0.0678 & \tblsecond{0.0687} & 0.0595 & 0.0658 & 0.0585 & -15.3\% \\
    \midrule
    \multirow{9}{*}{LastFM}
      & HR@5    & \tblbest{0.3638} & 0.2829 & 0.3522 & 0.3515 & 0.3526 & 0.3243 & 0.3097 & 0.3362 & 0.3430 & \tblsecond{0.3562} & -2.1\% \\
      & HR@10   & \tblbest{0.3994} & 0.3090 & 0.3780 & 0.3765 & \tblsecond{0.3831} & 0.3453 & 0.3312 & 0.3617 & 0.3657 & 0.3794 & -5.0\% \\
      & HR@20   & \tblbest{0.4252} & 0.3341 & 0.4038 & 0.3962 & \tblsecond{0.4089} & 0.3653 & 0.3508 & 0.3863 & 0.3909 & 0.4070 & -4.3\% \\
      & NDCG@5  & \tblsecond{0.3152} & 0.2589 & 0.3112 & 0.3086 & 0.2984 & 0.2956 & 0.2832 & 0.3067 & 0.3097 & \tblbest{0.3178} & +0.8\% \\
      & NDCG@10 & \tblbest{0.3267} & 0.2673 & 0.3196 & 0.3167 & 0.3083 & 0.3024 & 0.2900 & 0.3149 & 0.3171 & \tblsecond{0.3253} & -0.5\% \\
      & NDCG@20 & \tblbest{0.3332} & 0.2736 & 0.3260 & 0.3216 & 0.3149 & 0.3075 & 0.2950 & 0.3212 & 0.3234 & \tblsecond{0.3323} & -0.3\% \\
      & MRR@5   & \tblsecond{0.2988} & 0.2509 & 0.2976 & 0.2942 & 0.2803 & 0.2860 & 0.2742 & 0.2969 & 0.2986 & \tblbest{0.3049} & +2.0\% \\
      & MRR@10  & \tblsecond{0.3037} & 0.2543 & 0.3011 & 0.2976 & 0.2844 & 0.2888 & 0.2770 & 0.3002 & 0.3016 & \tblbest{0.3079} & +1.4\% \\
      & MRR@20  & \tblsecond{0.3054} & 0.2561 & 0.3027 & 0.2989 & 0.2863 & 0.2902 & 0.2784 & 0.3020 & 0.3033 & \tblbest{0.3099} & +1.5\% \\
    \midrule
    \multirow{9}{*}{KuaiRec}
      & HR@5    & 0.0945 & 0.0842 & \tblsecond{0.0951} & 0.0951 & 0.0940 & 0.0902 & 0.0902 & 0.0891 & 0.0907 & \tblbest{0.1000} & +5.2\% \\
      & HR@10   & \tblsecond{0.1055} & 0.0945 & 0.1011 & 0.1005 & 0.0995 & 0.0918 & 0.1011 & 0.0918 & 0.0967 & \tblbest{0.1060} & +0.5\% \\
      & HR@20   & \tblbest{0.1235} & 0.1109 & 0.1109 & 0.1158 & 0.1109 & 0.0923 & \tblsecond{0.1186} & 0.0995 & 0.1109 & \tblbest{0.1235} & +0.0\% \\
      & NDCG@5  & \tblsecond{0.0909} & 0.0756 & 0.0897 & 0.0902 & 0.0902 & 0.0885 & 0.0881 & 0.0863 & 0.0845 & \tblbest{0.0946} & +4.1\% \\
      & NDCG@10 & \tblsecond{0.0943} & 0.0790 & 0.0916 & 0.0919 & 0.0919 & 0.0890 & 0.0915 & 0.0871 & 0.0865 & \tblbest{0.0965} & +2.3\% \\
      & NDCG@20 & \tblsecond{0.0988} & 0.0831 & 0.0941 & 0.0958 & 0.0948 & 0.0892 & 0.0959 & 0.0891 & 0.0901 & \tblbest{0.1009} & +2.1\% \\
      & MRR@5   & \tblsecond{0.0897} & 0.0728 & 0.0879 & 0.0886 & 0.0890 & 0.0879 & 0.0874 & 0.0853 & 0.0824 & \tblbest{0.0928} & +3.5\% \\
      & MRR@10  & \tblsecond{0.0910} & 0.0742 & 0.0887 & 0.0893 & 0.0897 & 0.0882 & 0.0888 & 0.0857 & 0.0833 & \tblbest{0.0936} & +2.8\% \\
      & MRR@20  & \tblsecond{0.0922} & 0.0753 & 0.0894 & 0.0903 & 0.0905 & 0.0882 & 0.0899 & 0.0863 & 0.0843 & \tblbest{0.0948} & +2.7\% \\
    \bottomrule
  \end{tabular}}
  \caption{Full experimental results across all cutoffs. Bold and underlined values denote the best and one second-best result within each row.}
  \label{tab:appendix-full-seen}
\end{table*}

The appendix table makes it easier to verify whether improvements survive beyond the compact main-text summary without changing the selection rule between the main text and appendix.

\section{Extended Stage and Layout Ablations}
\label{app:stage-layout}
Figure~\ref{fig:appendix-stage-layout} expands the structural study in Figure~\ref{fig:stage-structure-panels} with the broader sweep of stage-removal and router-composition variants.

\begin{figure}[t]
  \centering
  \begin{minipage}[t]{0.48\linewidth}
    \centering
    \placeholderpanel{1.05in}

    \vspace{0.35em}
    {\footnotesize (a) Macro/mid/micro removal}
  \end{minipage}\hfill
  \begin{minipage}[t]{0.48\linewidth}
    \centering
    \placeholderpanel{1.05in}

    \vspace{0.35em}
    {\footnotesize (b) Dense FFN vs. 1/2/3-stage layouts}

  \end{minipage}

  \vspace{0.75em}

  \begin{minipage}[t]{\linewidth}
    \centering
    \placeholderpanel{1.15in}

    \vspace{0.35em}
    {\footnotesize (c) Router-composition and stage-order variants}
  \end{minipage}
  \caption{Extended stage and layout ablations. (a) Stage removal. (b) Dense replacement versus staged routing. (c) Router-composition and stage-order variants.}
  \Description{A three-panel placeholder figure for architecture ablations covering stage removal, dense-vs-stage-count comparisons, and router-composition or stage-order variants.}
  \label{fig:appendix-stage-layout}
\end{figure}

The figure should be read from left to right. Panel (a) verifies that macro, mid, and micro routing each contribute a distinct part of the gain. Panel (b) separates staged routing from generic expert insertion, and panel (c) checks that the final router composition is not an arbitrary configuration choice.

\section{Objective and Regularization Variants}
\label{app:objective}
Table~\ref{tab:appendix-objective} isolates the contribution of consistency, z-loss, and optional balance regularization using ranking quality, route consistency, and stability indicators. The main paper configuration corresponds to the consistency+z-loss row.

\begin{table}[t]
  \centering
  \caption{Objective and regularization variants.}
  \label{tab:appendix-objective}
  \small
  \setlength{\tabcolsep}{5pt}
  \begin{tabularx}{\columnwidth}{@{}>{\raggedright\arraybackslash}XYYY@{}}
    \toprule
    Variant & MRR@20 & Consistency & Stability \\
    \midrule
    No auxiliary loss & -- & -- & -- \\
    Consistency only & -- & -- & -- \\
    z-loss only & -- & -- & -- \\
    Balance loss only & -- & -- & -- \\
    Consistency + z-loss (main configuration) & -- & -- & -- \\
    Consistency + z-loss + balance & -- & -- & -- \\
    \bottomrule
  \end{tabularx}
\end{table}

This compact table shows whether the selected main objective improves MRR@20 while preserving routing consistency and stability, and whether adding explicit balance pressure helps or instead weakens behavior-aligned specialization.

\section{Extended Diagnostics}
\label{app:diagnostics}
Figure~\ref{fig:appendix-diagnostics} groups the main routing diagnostics into one coordinated layout so that stage usage, sharpness, and semantic consistency can be read together.

\begin{figure}[t]
  \centering
  \begin{minipage}[t]{0.48\linewidth}
    \centering
    \placeholderpanel{1.18in}

    \vspace{0.35em}
    {\footnotesize (a) Expert-usage heatmaps by stage}
  \end{minipage}\hfill
  \begin{minipage}[t]{0.48\linewidth}
    \centering
    \placeholderpanel{1.18in}

    \vspace{0.35em}
    {\footnotesize (b) Entropy and effective-expert-count summary}
  \end{minipage}

  \vspace{0.8em}

  \begin{minipage}[t]{0.48\linewidth}
    \centering
    \placeholderpanel{1.18in}

    \vspace{0.35em}
    {\footnotesize (c) Stage-wise route consistency}
  \end{minipage}\hfill
  \begin{minipage}[t]{0.48\linewidth}
    \centering
    \placeholderpanel{1.18in}

    \vspace{0.35em}
    {\footnotesize (d) Feature-bucket route patterns}
  \end{minipage}
  \caption{Routing diagnostics. (a) Expert usage by stage. (b) Entropy and effective expert count. (c) Stage-wise route consistency. (d) Feature-bucket routing patterns.}
  \Description{A four-panel placeholder figure for expert-usage heatmaps, entropy summaries, stage-wise consistency, and feature-bucket route patterns.}
  \label{fig:appendix-diagnostics}
\end{figure}

These diagnostics are organized from coarse structure to semantic interpretation. The first two panels show how broad or sharp routing is at each stage, while the last two panels connect those patterns back to feature similarity and interpretable behavioral regimes.

\section{Behavior-Slice Analysis}
\label{app:behavior-slices}
Because RouteRec is motivated by heterogeneous behavior, a slice-wise view is a natural complement to the aggregate metrics. This appendix section is intended to localize gains to behavior regimes such as repeat-heavy prefixes, high-tempo sessions, narrow-focus sessions, and exploration-heavy prefixes.

\begin{figure}[t]
  \centering
  \begin{minipage}[t]{\linewidth}
    \centering
    \placeholderpanel{1.25in}

    \vspace{0.35em}
    {\footnotesize (a) Slice-wise MRR@20 across behavioral regimes}
  \end{minipage}\hfill

  \vspace{0.75em}

  \begin{minipage}[t]{\linewidth}
    \centering
    \placeholderpanel{1.25in}

    \vspace{0.35em}
    {\footnotesize (b) Relative RouteRec gain and route concentration by slice}
  \end{minipage}
  \caption{Behavior-slice analysis. (a) Slice-wise ranking quality. (b) Relative gain and routing concentration by slice.}
  \Description{A two-panel placeholder figure for slice-wise MRR@20 and RouteRec gain across repeat-heavy, fast-tempo, focused, and exploration-heavy regimes.}
  \label{fig:appendix-behavior-slices}
\end{figure}

The useful reading order is to first check which slices benefit most in panel (a), then verify in panel (b) whether those gains coincide with more decisive and semantically aligned routing. This experiment fits the paper's main argument especially well because it answers not only whether RouteRec helps, but also where the help appears.

\section{Extended Transfer Results}
\label{app:transfer}
This section reports additional transfer experiments for low-resource and cross-domain adaptation.

\begin{table}[t]
  \centering
  \caption{Source-target transfer matrix.}
  \label{tab:appendix-transfer-matrix}
  \scriptsize
  \setlength{\tabcolsep}{3pt}
  \begin{tabular*}{\columnwidth}{@{\extracolsep{\fill}}lcccccc@{}}
    \toprule
    Src. $\rightarrow$ Tgt. & Bty & Fsq & KRec & LFM & ML1M & RR \\
    \midrule
    Bty & -- & -- & -- & -- & -- & -- \\
    Fsq & -- & -- & -- & -- & -- & -- \\
    KRec & -- & -- & -- & -- & -- & -- \\
    LFM & -- & -- & -- & -- & -- & -- \\
    ML1M & -- & -- & -- & -- & -- & -- \\
    RR & -- & -- & -- & -- & -- & -- \\
    \bottomrule
  \end{tabular*}
\end{table}

The matrix provides the exact source-target view that is too detailed for the main paper. It is useful for identifying asymmetric transfer pairs and for checking whether router reuse is most effective between domains with similar temporal and repetition structure. For compactness, the table uses the abbreviations Bty, Fsq, KRec, LFM, ML1M, and RR.

\begin{figure}[t]
  \centering
  \begin{minipage}[t]{0.48\linewidth}
    \centering
    \placeholderpanel{1.05in}

    \vspace{0.35em}
    {\footnotesize (a) Low-resource transfer curves}
  \end{minipage}\hfill
  \begin{minipage}[t]{0.48\linewidth}
    \centering
    \placeholderpanel{1.05in}

    \vspace{0.35em}
    {\footnotesize (b) Frozen/finetuned/anchor and group/full router transfer}

  \end{minipage}

  \vspace{0.75em}

  \begin{minipage}[t]{\linewidth}
    \centering
    \placeholderpanel{1.15in}

    \vspace{0.35em}
    {\footnotesize (c) Multi-source sequential order effect}
  \end{minipage}
  \caption{Transfer variants. (a) Low-resource transfer curves. (b) Frozen, finetuned, anchor, and router-granularity comparisons. (c) Multi-source order effects.}
  \Description{A three-panel placeholder figure for low-resource transfer curves, transfer-strategy comparisons, and multi-source order effects.}
  \label{fig:appendix-transfer-panels}
\end{figure}

Together, the table and figure let the reader move from exact pairwise results to higher-level patterns. The intended appendix message is that transfer is asymmetric, most helpful in related domains and low-resource regimes, and partially preserved even under coarser router reuse.

\end{document}
